\section{Related work and remaining obligations}
\label{sec:related}

The general hinged-dissection landscape is much broader than this paper.  Hinged
dissections exist in substantial generality~\cite{abbott2007hinged}; reversible
hinged dissections and reversible nets are studied in
\cite{akiyama2018polyhedral,akiyama2016reversible}; and standard background on
geometric folding, linkages, and polyhedra is given by Demaine and
O'Rourke~\cite{demaine2007gfalop}.  Continuous flattening, vertex-unfolding,
and digital jointed assembly systems provide adjacent contexts
\cite{demaine2024continuous,demaine2001vertex,rodsteward2019,joinable2021}.

The novelty candidate here is not a new universal hinged-dissection theorem.
It is the combination of a fixed tetrahedral median-plane model, exact
zero-thickness rooted hinge-tree kinematics, row-level route-wrapper
certificates, and a public companion repository that keeps theorem evidence separate
from digital fabrication artifacts.

\paragraph{Remaining obligations.}
The following are deliberately open or out of scope.
\begin{itemize}[leftmargin=*]
  \item Positive-thickness hinge design and physical fabrication validation.
  \item Global collision-free motion beyond the audited one-parameter domain.
  \item Dynamic connectedness between \treeA{} and \treeB{}.
  \item Three-parameter bounded-cell closure.
  \item Extension of the B04 bridge beyond the six selected rows.
  \item A general theorem for all S4 hinge trees or all tetrahedral dissections.
\end{itemize}

The package is therefore best read as a scoped certificate paper: it identifies
what has been proved for a concrete object, ships the evidence in a reproducible
layout, and makes the boundaries visible enough for independent review.
